%% IntSeqBERT — CICM 2026 presentation (Beamer / metropolis, 16:9)
%% Compile: pdflatex slides && pdflatex slides

\documentclass[aspectratio=169,10pt]{beamer}

\usetheme[progressbar=frametitle,numbering=fraction]{metropolis}

\usepackage[T1]{fontenc}
\usepackage{lmodern}
\usepackage{amsmath}
\usepackage{amssymb}
\usepackage{booktabs}
\usepackage{graphicx}
\usepackage{pgfpages}

% Keep the 16:9 slide layout, but emit each page on A4 landscape paper.
% This also makes printer-driver n-up layouts (such as 4 pages per sheet)
% handle the document at the intended physical size.
\pgfpagesuselayout{resize to}[a4paper,landscape]

\graphicspath{{../figures/}}

% slightly tighter itemize spacing for text-rich slides
\newcommand{\tightlist}{\setlength{\itemsep}{0.2em}}

\title{IntSeqBERT: Learning Arithmetic Structure in OEIS\texorpdfstring{\\}{ }via Modulo-Spectrum Embeddings}
\author{Kazuhisa Nakasho}
\institute{Iwate Prefectural University, Japan\\[1ex]
  Code \& data: \texttt{github.com/aabaa/IntSeqBERT}}
\date{CICM 2026}

\begin{document}

\begin{frame}
  \titlepage
\end{frame}

\begin{frame}{Why Finding Good Problems Matters}
  \small
  \textbf{Terence Tao's concern:}
  We may solve fruitful open problems faster than we develop new ones.
  \begin{quote}
    ``a country or region can suffer a critical shortage of drinking water
    while simultaneously being surrounded by a massive ocean.''
  \end{quote}
  \begin{itemize}\tightlist
    \item We can ask \textbf{infinitely many mathematical questions}.
    \item But \alert{good problems may still become scarce}:
          many questions lead to few new ideas or connections.
    \item A useful problem should help us \textbf{understand more mathematics}.
  \end{itemize}
  \begin{alertblock}{Our motivation}
    Can AI help us find new mathematical connections and promising conjectures?
  \end{alertblock}
  {\scriptsize Source: Terence Tao, Mastodon, post 1/4. Quote above; other text is a summary.\\
  \href{https://mathstodon.xyz/@tao/117237320796901560}{\texttt{mathstodon.xyz/@tao/117237320796901560}}}
\end{frame}

\begin{frame}{Long-Term Goal --- AI for Mathematical Discovery}
  \small
  \begin{itemize}\tightlist
    \item Many important conjectures began with a \alert{numerical observation}:
      \begin{itemize}
        \item $196{,}884 = 196{,}883 + 1$ --- McKay's observation $\to$ \emph{Monstrous Moonshine}
        \item The BSD conjecture and mirror symmetry ($317{,}206{,}375$ rational curves)
              also began with numerical patterns
      \end{itemize}
    \item Today, AI systems take part in this kind of discovery:\\
          \textbf{AlphaProof}, the \textbf{Ramanujan Machine}, \textbf{FunSearch}
    \item \alert{Our long-term goal}: an AI that finds \emph{hidden number-theoretic
          similarities} between integer sequences from different fields ---
          and supports \alert{conjecture generation}
    \item Our data source is the \textbf{OEIS}: 391{,}710 sequences, each with a precise
          definition --- a collection of mathematics that computers can read
    \item We need a machine \alert{representation} that captures
          number-theoretic structure
    \item \textbf{This paper is the first step}: we propose such a representation,
          and we \emph{measure} how much arithmetic structure it actually captures
  \end{itemize}
\end{frame}

\begin{frame}{Why Multiplication Is Hard for Deep Learning}
  \footnotesize
  \begin{itemize}\tightlist
    \item The main problem: standard neural models have difficulty learning
          \alert{multiplication-like rules}.
    \item \textbf{Spectral bias}: networks learn low-frequency, addition-like structure
          first (Rahaman+ 2019)
    \item How many layers does multiplication need?
      \begin{itemize}
        \item \textbf{In theory}: \alert{depth-3 threshold circuits} are enough for
              $n$-bit multiplication, and fewer layers are not enough (Siu \& Roychowdhury 1994)
        \item \textbf{In practice}: gradient descent does not find these circuits ---
              GPT-4 scores 59\% on $3{\times}3$-digit, and only \alert{4\% on
              $4{\times}4$-digit} multiplication (Dziri+ 2023)
        \item Even \textbf{mod-$p$ multiplication in a 2-layer Transformer} generalises only
              after a very long training period --- the \emph{grokking} phenomenon (Power+ 2022)
        \item One successful method uses \textbf{12 layers \emph{plus} a special training method}
              (implicit chain-of-thought) to reach 15-digit products (Yao+ 2025)
      \end{itemize}
    \item \textbf{Tokenisation makes it worse}: values outside the vocabulary become
          \texttt{UNK}, and token IDs do not directly show arithmetic structure
    \item The OEIS is a \alert{difficult test}: its values range from
          single digits \dots\ to $10^{210}$
  \end{itemize}
\end{frame}

\begin{frame}{Key Idea --- Show Multiplicative Structure in the Input}
  \begin{itemize}\tightlist
    \item Our key idea comes from \alert{classical number theory}.
    \item Reduction mod $m$ is a \alert{ring homomorphism}
          $\mathbb{Z} \to \mathbb{Z}/m\mathbb{Z}$:
          it preserves \textbf{both addition and multiplication}
    \item So the residue sequence $(x_i \bmod m)$ keeps the multiplicative structure ---
          \textbf{even for very large values}: for $10^{210}$, the residue stays in
          $\{0,\dots,m-1\}$
    \item Number theory explains these \textbf{periodic patterns}:
      \begin{itemize}
        \item \textbf{Fermat's little theorem}: powers are periodic mod $p$
              (period divides $p-1$)
        \item \textbf{Quadratic reciprocity}: links residue information across primes
        \item \textbf{CRT}: composite moduli combine information from several prime powers
      \end{itemize}
    \item \alert{Hypothesis}: give the network the \emph{modulo spectrum} as
          input --- then it sees multiplicative structure \alert{at depth zero},
          without having to learn it across many layers
  \end{itemize}
\end{frame}

\begin{frame}{Task --- Masked Sequence Modelling on OEIS}
  \begin{itemize}\tightlist
    \item The task is \alert{masked sequence modelling} --- BERT-style, but for integers.
    \item We take up to the first 128 terms of an OEIS sequence, mask 15\% of the
          positions, and train the model to predict each masked value from context
    \item This tests how well the model has learned the \emph{rule} of the
          sequence; \textbf{next-term prediction} is the special case where the last term
          is masked
    \item Prior benchmark: \textbf{FACT} (NeurIPS 2022) --- a plain Transformer on five OEIS
          tasks; \emph{unmasking} was the hardest; our \textbf{Vanilla baseline}
          uses that architecture
    \item Instead of predicting one token, we predict \alert{three quantities} per masked
          position:
      \begin{enumerate}
        \item \textbf{Magnitude}: $v_i = 1 + \log_{10}|x_i|$
        \item \textbf{Sign}: $s_i \in \{+, -, 0\}$
        \item \textbf{Residues}: $x_i \bmod m$ for every $m \in \{2, \dots, 101\}$
      \end{enumerate}
  \end{itemize}
\end{frame}

\begin{frame}{IntSeqBERT --- Dual-Stream Architecture}
  \centering
  \includegraphics[width=0.8\textwidth]{fig1_architecture.png}
  \begin{itemize}\tightlist\footnotesize
    \item The \alert{magnitude stream} ($\mathbb{R}^4$) encodes log-scale size and sign
          $\to$ MLP
    \item The \alert{modulo stream} ($\mathbb{R}^{200}$) embeds each residue on the
          \textbf{unit circle}: $\phi_m(r) = [\sin(2\pi r/m), \cos(2\pi r/m)]$
          --- we use the Transformer's sinusoidal \emph{positional} encoding for
          \emph{residues of values}; residues $m{-}1$ and $0$ stay neighbours
    \item \alert{FiLM fusion}: the modulo stream scales and shifts the magnitude stream:
          $\mathbf{e}_i = (1 + \boldsymbol{\gamma}_i) \odot \mathbf{h}_i^{\mathrm{mag}}
           + \boldsymbol{\beta}_i$
    \item Then a standard \textbf{pre-LN Transformer encoder} ---
          Small 6.4M / Middle 29.0M / \textbf{Large 91.5M} parameters
  \end{itemize}
\end{frame}

\begin{frame}{Training --- Multi-Task Objective}
  \begin{itemize}\tightlist
    \item Three prediction heads use the encoder output:
      \begin{itemize}
        \item The \textbf{magnitude head} also predicts its own \alert{uncertainty}
              $(\mu_i, \sigma_i^2)$ --- remember this: the Solver will use it
        \item The \textbf{sign head}: 3-class classification
        \item The \textbf{modulo head}: 100 independent classifiers, each normalised by its
              maximum entropy $\ln m$ --- so mod 2 and mod 101 contribute fairly
      \end{itemize}
    \item The loss is a weighted sum:
          $\mathcal{L} = \mathcal{L}_{\mathrm{mag}} + \mathcal{L}_{\mathrm{sign}}
           + 2\,\mathcal{L}_{\mathrm{mod}}$
    \item Trained on \textbf{219{,}765 OEIS sequences}, 200 epochs, in \textbf{FP32} ---
          because values near $10^{210}$ overflow FP16
    \item All experiments run on a \alert{single consumer GPU}
          (RTX 3070\,Ti, 8\,GB), so other researchers can repeat the experiments
  \end{itemize}
\end{frame}

\begin{frame}{Neural-Symbolic Decoding --- the CRT Solver}
  \begin{itemize}\tightlist
    \item The network outputs \textbf{distributions} --- magnitude, sign, 100 residues ---
          \alert{not integers}
    \item The \alert{Solver} produces integer candidates using \textbf{symbolic rules}:
      \begin{itemize}
        \item The magnitude and its uncertainty give a $3\sigma$ \textbf{search interval}
        \item \textbf{Dense} ($\Delta n \le 10^6$): check every candidate in the interval
        \item \textbf{Sieve} ($\le 10^{14}$): CRT beam search based on the moduli with the
              most confident predictions
        \item \textbf{CRT} ($> 10^{14}$): sparse CRT generates very large candidates directly
      \end{itemize}
    \item We combine neural prediction and symbolic rules in a \alert{neural-symbolic} system ---
          residue arithmetic and the Chinese Remainder Theorem are \textbf{built in};
          the network does not have to learn these rules from data
    \item Is this a strength or a limitation? ---
          \alert{I will come back to this with the results.}
  \end{itemize}
\end{frame}

\begin{frame}{Dataset --- 274{,}705 OEIS Sequences}
  \begin{itemize}\tightlist
    \item From 391{,}710 OEIS entries we keep \alert{274{,}705 (70.1\%)}, split 8:1:1
    \item Eleven tags are excluded, in \textbf{three groups}:
      \begin{itemize}
        \item \textbf{depend on how values are written} (\texttt{cons}, \texttt{cofr},
              \texttt{frac}, \texttt{base}, \texttt{word}):
              their patterns depend on the \emph{encoding} of the values
        \item \textbf{do not fit this task} (\texttt{fini}, \texttt{tabl})
        \item \textbf{quality labels} (\texttt{dead}, \texttt{unkn}, \texttt{less},
              \texttt{dumb})
      \end{itemize}
    \item These choices define our task:
          \texttt{less}/\texttt{dumb} are subjective labels, and \texttt{tabl} could be
          converted into sequences. Future benchmarks may use different choices
    \item For evaluation, five \textbf{magnitude buckets} (groups by size):\\
          Small ($<10^2$, 52.7\%) $\to$ Medium $\to$ Large $\to$ Huge $\to$
          Astronomical ($\ge 10^{50}$)
  \end{itemize}
\end{frame}

\begin{frame}{Metrics --- and Three Questions}
  \begin{itemize}\tightlist
    \item \alert{Magnitude Accuracy ($\mathrm{Acc}_{0.5}$) is NOT exact match}:
          it asks whether the prediction is within $0.5$ on the log scale ---
          a factor of $\sqrt{10} \approx 3.16$
      \begin{itemize}
        \item Example: true value $10^6$ $\to$ anything in
              $[3.2 \times 10^5,\ 3.2 \times 10^6]$ counts
        \item Its purpose: help set a \textbf{realistic search range} for the Solver
      \end{itemize}
    \item \textbf{Exact} integer match is measured separately, by \textbf{Solver Top-$k$}
    \item \textbf{MMA}: mean classification accuracy over all 100 moduli
    \item \textbf{NIG} (normalised information gain)
          $= 1 - \mathcal{L}_{\mathrm{CE}}^{(m)}/\ln m$:
          0 = no gain over uniform prediction, 1 = perfect prediction
  \end{itemize}
  \begin{alertblock}{Three questions for the results}
    \textbf{Q1} --- Does the dual stream help?\quad
    \textbf{Q2} --- Where does it help?\quad
    \textbf{Q3} --- Can we predict exact integers?
  \end{alertblock}
\end{frame}

\begin{frame}{Q1 --- Does It Help? Yes, Consistently}
  \begin{center}\small
    \begin{tabular}{llrrr}
      \toprule
      Size & Model & Mag Acc & Sign Acc & MMA \\
      \midrule
      Large & \textbf{IntSeqBERT} & \textbf{95.85} & \textbf{98.54} & \textbf{50.38} \\
      Large & Vanilla             & 86.97          & 97.66          & 45.85 \\
      Large & Ablation (no modulo)& 89.70          & 98.29          & 35.22 \\
      \bottomrule
    \end{tabular}
  \end{center}
  \begin{itemize}\tightlist
    \item At the Large scale, IntSeqBERT reaches \alert{95.85\%} magnitude accuracy ---
          \textbf{+8.9\,pt} over the tokenised baseline ---
          and \alert{50.38\%} mean modulo accuracy (\textbf{+4.5\,pt})
    \item The \alert{ablation row} is the important one: remove the modulo stream $\to$
          MMA drops \alert{$-15.2$\,pt} --- and magnitude accuracy \textbf{also} drops
          \textbf{$-6.2$\,pt}
    \item The residues limit the possible integer values. This helps the model learn
          \alert{both magnitude and modulo information}
    \item (The same pattern holds at Small and Middle scale --- scaling results in backup)
  \end{itemize}
\end{frame}

\begin{frame}{Q2 --- Where? Exactly Where Tokenisation Breaks}
  \centering
  \includegraphics[width=0.78\textwidth]{fig5_magnitude_scatter.pdf}
  \begin{columns}[T]
    \begin{column}{0.34\textwidth}
      \centering\scriptsize
      \begin{tabular}{lrrr}
        \toprule
        MSE & IntSeq & Vanilla & Abl. \\
        \midrule
        Small & 0.111          & 0.138 & \textbf{0.103} \\
        Large & \textbf{0.162} & 2.100 & 0.381 \\
        Huge  & \textbf{2.082} & 22.73 & 5.021 \\
        \bottomrule
      \end{tabular}
    \end{column}
    \begin{column}{0.62\textwidth}
      \scriptsize
      \begin{itemize}\setlength{\itemsep}{0.2em}
        \item Above its vocabulary limit, Vanilla uses \texttt{UNK} ---
              its Large-bucket error is \alert{13$\times$ ours}; overall
              $R^2$: \alert{0.988} vs 0.943
        \item This is why bucketed evaluation matters: most OEIS sequences grow slowly,
              so easier cases have a large effect on the \emph{overall} scores
        \item The clearest gains are in the \alert{hard buckets --- Large and
              above} --- where the baseline stops working
      \end{itemize}
    \end{column}
  \end{columns}
\end{frame}

\begin{frame}{Q3 --- The Strictest Test: Exact Next-Term Prediction}
  \small
  \begin{center}\small
    \begin{tabular}{llrr}
      \toprule
      Size & Model & Top-1 & Top-10 \\
      \midrule
      Large & \textbf{IntSeqBERT} & \textbf{19.09} & \textbf{26.23} \\
      Large & Vanilla             & 2.59           & 3.80 \\
      Large & Ablation            & 11.75          & 21.79 \\
      \bottomrule
    \end{tabular}
  \end{center}
  \begin{itemize}\tightlist
    \item We test the full system: predict the \textbf{exact last term} of 10{,}000
          sequences kept separate for testing
    \item The tokenised Transformer gets \textbf{2.59\%}; IntSeqBERT gets \textbf{19.09\%}
          --- a \alert{7.4-fold improvement}
    \item By bucket: Small \textbf{68.3\%}, Medium \textbf{20.8\%} ---
          while Vanilla is at \textbf{0\%} for everything beyond its vocabulary
    \item But for the Large buckets and above, our model also drops to near zero:
          CRT reconstruction needs \textbf{exact residues for many moduli at the same time}
          --- and at MMA $\approx 50\%$, that almost never happens
    \item The next question: \alert{is 19\% a success or a failure?}
  \end{itemize}
\end{frame}

\begin{frame}{What Does 19\% Mean? --- A Useful Negative Result}
  \begin{itemize}\tightlist
    \item There are \alert{two useful ways to understand this result.}
    \item Even with number-theoretic structure \textbf{built into} the representation,
          exact prediction of OEIS terms \textbf{remains hard};
          a purely tokenised Transformer performs poorly (2.59\%)
    \item \textbf{Point 1}: the results measure \alert{how limited end-to-end learning is}
          for arithmetic in this setting
    \item \textbf{Point 2}: the results support \alert{neural-symbolic design} ---
          the CRT solver provides exact rules that the network
          does not learn well on its own
    \item And my answer to the earlier question --- strength or limitation?\\
          We use features based on number theory as a \alert{tool for measurement}:
          removing the modulo stream shows how much these features help
          the model learn arithmetic structure ($-15.2$\,pt MMA)
  \end{itemize}
\end{frame}

\begin{frame}{But What Did the Network Learn? The NIG Spectrum}
  \centering
  \includegraphics[width=0.72\textwidth]{fig2_nig_spectrum.pdf}
  \begin{itemize}\tightlist\footnotesize
    \item Prediction accuracy is one way to evaluate the model. We also ask:
          \alert{what did the network actually learn about number theory?}
    \item NIG for every modulus $m = 2, \dots, 101$: IntSeqBERT is above both baselines
          across the \textbf{entire spectrum} --- and the shape is not flat:
          \alert{the peaks follow a clear pattern}
    \item Best modulus across \emph{all} models and scales:
          $m = 96 = 2^5 \times 3$ (NIG\,$=0.629$); second best: $m = 2$ ---
          \textbf{parity is everywhere} in the OEIS
    \item Parity also shows differences between the models: mod-2 accuracy
          \textbf{85.7} (IntSeq) / 81.4 (Vanilla) / 72.1 (Ablation)
  \end{itemize}
\end{frame}

\begin{frame}{Finding --- Higher NIG with a Lower Totient Ratio}
  \begin{columns}[T]
    \begin{column}{0.46\textwidth}
      \centering
      \includegraphics[width=\textwidth]{fig2b_nig_vs_phi.pdf}
    \end{column}
    \begin{column}{0.54\textwidth}
      \footnotesize
      \begin{itemize}\setlength{\itemsep}{0.4em}
        \item Why 96? Why composites?
              \alert{One quantity helps explain this pattern}:
              Euler's totient ratio $\varphi(m)/m = \prod_{p \mid m}(1 - 1/p)$
        \item NIG vs $\varphi(m)/m$: Pearson \alert{$r = -0.851$} ($p < 10^{-28}$)
        \item We explain this using the \alert{CRT aggregation effect}:
              $x \bmod 96$ gives us
              $x \bmod 2, 4, 8, 16, 32$ --- \emph{and} $x \bmod 3$ --- at the same time
        \item This is an experimental result, not a theorem. It suggests a
              \textbf{design choice}: to extend beyond 101, add
              \alert{highly composite numbers}, not more primes
        \item This spectrum describes the \alert{arithmetic patterns} in a dataset.
              We can use spectra to \textbf{compare sequences or datasets}
      \end{itemize}
    \end{column}
  \end{columns}
\end{frame}

\begin{frame}{Positioning in Computer Mathematics}
  \small
  \begin{itemize}\tightlist
    \item Comparing sequences brings us back to our original goal:
          \alert{what is this work for?}
    \item \textbf{Program synthesis on the OEIS} --- Alien Coding, QSynt, Learning
          Conjecturing --- describes each sequence's \textbf{generation rule}
    \item Our work can \alert{support these methods}: we learn \alert{shared arithmetic structure}
          across the whole corpus
      \begin{itemize}
        \item Synthesis asks: \emph{``What program generates this sequence?''}
        \item Our representation asks: \emph{``Which sequences are number-theoretically
              \textbf{similar} --- even if generated in completely different ways?''}
      \end{itemize}
    \item Pretrained embeddings could \textbf{guide the search} in program synthesis,
          or help us \alert{find similar sequences}
    \item Finding similar sequences could help us \alert{generate conjectures} ---
          our long-term goal
    \item Code, preprocessing, and evaluation: \texttt{github.com/aabaa/IntSeqBERT}
  \end{itemize}
\end{frame}

\begin{frame}{Conclusion}
  \small
  \textbf{Four main points:}
  \begin{enumerate}\tightlist
    \item The modulo spectrum uses features \alert{based on number theory}:
          a \textbf{ring homomorphism} makes multiplicative structure visible
          \alert{at depth zero}
    \item It works: consistent gains at every scale, and a \alert{7.4$\times$}
          improvement in exact next-term prediction
    \item \textbf{Composite moduli combine arithmetic information}:
          NIG has a negative correlation with $\varphi(m)/m$ ($r = -0.851$)
    \item At the same time: pure neural prediction of OEIS terms \textbf{remains hard} ---
          this supports further work on \alert{neural-symbolic methods}
  \end{enumerate}
  \medskip
  \textbf{Future work:} approximate CRT for large integers \,$\cdot$\,
  highly composite moduli beyond 101 \,$\cdot$\,
  larger-scale retraining \,$\cdot$\,
  \alert{similarity retrieval \& conjecture generation}
  \bigskip
  \begin{center}
    \alert{\textbf{Thank you!}} --- \texttt{github.com/aabaa/IntSeqBERT}
  \end{center}
\end{frame}

\appendix

\begin{frame}{Backup --- Full Results \& Scaling Behaviour}
  \centering
  \includegraphics[width=0.75\textwidth]{fig3_scaling.pdf}
  \begin{itemize}\tightlist\footnotesize
    \item The same ranking holds at every scale: IntSeqBERT $>$ Vanilla $>$ Ablation
    \item \alert{Model size helps modulo prediction more}: from Small to Large, MMA improves
          \textbf{+10.0\,pt} while Mag Acc improves only \textbf{+1.1\,pt} ---
          larger models improve modulo prediction more than magnitude prediction
    \item Solver Top-1 also improves: \textbf{+5.0\,pt} (14.05 $\to$ 19.09)
    \item No clear sign of overfitting: Train 1.00 / Val 1.01 at epoch 200;
          larger-scale retraining is planned as future work
  \end{itemize}
\end{frame}

\begin{frame}{Backup --- Why Not Digit-Based Tokenisation?}
  \begin{itemize}\tightlist
    \item A natural alternative: treat each decimal digit as a token, and generate
          digit by digit
    \item We have not made an experimental comparison in this paper --- it is future work
    \item Our target is not the pattern of decimal digits, but the
          \alert{number-theoretic structure of the integer itself}:
      \begin{itemize}
        \item the modulo stream supplies residues \textbf{directly};
              with digit tokens, the model must \textbf{calculate} them from the digits
      \end{itemize}
    \item Digit-by-digit generation also \textbf{changes the problem setting}:
          sequential prediction of a digit string vs.\ predicting one integer value
    \item We plan to compare the two settings under the same conditions
  \end{itemize}
\end{frame}

\begin{frame}{Backup --- Limitations}
  \begin{itemize}\tightlist
    \item \textbf{Large integers}: CRT reconstruction needs exact residues for many moduli
          at the same time; 13.4\% of Solver calls return no valid candidate
          (\texttt{none} mode); approximate CRT / less strict residue constraints are
          possible next steps
    \item \textbf{Dataset bias}: many \texttt{nonn} sequences $\to$ high sign accuracy;
          many slowly growing sequences $\to$ high overall Mag Acc
    \item \textbf{Compute}: a single 8\,GB GPU, single seed (42), vocabulary 20{,}003 ---
          Vanilla uses the \emph{same architecture as FACT} with our settings.
          We do not reproduce FACT's reported results
  \end{itemize}
\end{frame}

\begin{frame}{Backup --- How Hard Is Multiplication for Neural Networks?}
  \small
  \begin{itemize}\tightlist
    \item \textbf{Circuit theory}: $n$-bit multiplication is \textbf{not in AC$^0$}
          (constant-depth AND/OR), but \textbf{is} computable by polynomial-size
          \alert{depth-3 threshold circuits} --- and depth 3 is optimal
          (Siu \& Roychowdhury, \emph{SIAM J.\ Discrete Math.}, 1994);
          iterated products: depth 4
    \item \textbf{Transformers $\approx$ TC$^0$}: log-precision transformers are
          upper-bounded by uniform TC$^0$ (Merrill \& Sabharwal 2023) ---
          this connects multiplication to the limits of these models
    \item \textbf{LLM experiments} (Dziri+, \emph{Faith and Fate}, NeurIPS 2023):
          GPT-4 zero-shot --- $3{\times}3$-digit: \textbf{59\%},
          $4{\times}4$-digit: \alert{4\%}
    \item \textbf{Grokking} (Power+ 2022, arXiv:2201.02177): a 2-layer transformer on
          modular arithmetic generalises only after a very long training period
    \item \textbf{Mechanistic analysis} (arXiv:2510.00184, 2025): standard training fails
          at multi-digit multiplication even at 90B parameters; a 12-layer model succeeds
          at 15-digit products only with implicit-CoT training, forming a
          \alert{binary tree of partial products} in attention
  \end{itemize}
\end{frame}

\end{document}
